% Generated by tools/edn_latex.py from reports/edn.md.
% Do not edit: change reports/edn.md and run `make edn`.
\ednentry{
  id = {EDN-56},
  title = {The model bundle carries its own provenance, so \texorpdfstring{\texttt{models/\allowbreak{}}}{models/} is deliberately not byte-reproducible},
  date = {2026-09-30},
  milestone = {M3: Quality Assurance},
  topic = {Reproducibility, Data Versioning},
  participants = {@lukas2510, decided by the repository owner},
  decision = {\texttt{models/\allowbreak{}<variant>/\allowbreak{}model.\allowbreak{}json} carries \texttt{trained\_\allowbreak{}at} and \texttt{mlflow.\allowbreak{}run\_\allowbreak{}id} alongside the hyperparameters, the feature list and the training statistics. Both change between two otherwise identical fits, so the model directory is \textbf{not} byte-reproducible by design. The reproducibility claim NFR-06 rests on is made about the predictions and about the estimator's own payload file instead, and both are asserted.},
  rationale = {\emph{Alternatives considered:}
\begin{itemize}[leftmargin=*, topsep=2pt]
\item \textbf{Option A (chosen): provenance inside \texttt{model.\allowbreak{}json}.}
\newline Pros: the bundle is self-describing wherever it ends up. EDN-08 bakes it into the API image at build time, so at serving time there is no MLflow client, no \texttt{dvc.\allowbreak{}lock} and no repository -- and \texttt{GET /\allowbreak{}health} still has to be able to say which model version it serves and where its numbers came from. It is also what lets \texttt{evaluate} resume the right run (EDN-55) without a second artefact to keep in step.
\newline Cons: \texttt{dvc repro train@\allowbreak{}<variant>} always produces a new directory hash even when the model is identical, so \texttt{evaluate} always reruns after a retrain, and \texttt{dvc status} can never say ``the model is unchanged''. A reviewer comparing two pipeline states cannot tell a real change from a re-run by the hash alone.
\item \textbf{Option B: keep \texttt{models/\allowbreak{}} byte-stable and move the provenance to a second, non-cached output.}
\newline Pros: the model directory's hash then means exactly ``the model changed'', which is the property a reviewer and \texttt{dvc status} both want, and a re-run that changes nothing reruns nothing downstream.
\newline Cons: it needs a second \texttt{outs} entry per variant, and it puts the run id in a file \texttt{evaluate} and the API would each have to find beside the model rather than in it -- a cross-stage dependency for \#39 and a second file the API image has to be told to copy. It also splits one artefact that is restored together into two that might not be.
\item \textbf{Option C: drop the provenance and rely on MLflow for it.}
\newline Pros: \texttt{models/\allowbreak{}} becomes byte-stable at no structural cost, and MLflow already records the commit, the data version and the parameters of every run.
\newline Cons: it works only where MLflow is reachable. A bundle obtained with \texttt{dvc pull} on a machine with no credentials -- which is the normal case for this project, see EDN-54 -- could then not say which run produced it, and neither could the running API. It would also break the resumption of EDN-55 entirely, because the id would live only on the server whose runs it is meant to identify.
\end{itemize}
\par
\emph{Why:} This is a trade against NFR-06 and it is worth being exact about which half of NFR-06 it trades. NFR-06 asks for two things: that the documented steps from a clean clone reproduce the same splits and metrics within 0.1 percentage points, and that \textbf{every training run records which code, which data version and which parameters produced it}. The provenance fields serve the second half directly. What they cost is a property NFR-06 never asks for, which is that the output directory hash is stable.
\par
The reproducibility that is claimed is therefore claimed about the right things and is tested. Two fits of the same variant on the same rows give bit-identical predictions for all four variants, and a byte-identical payload file -- \texttt{booster.\allowbreak{}txt}, \texttt{pipeline.\allowbreak{}joblib} or \texttt{lookup.\allowbreak{}parquet}, which is the whole of the fitted model. \texttt{model.\allowbreak{}json} is excluded from that comparison explicitly, in the test and in its docstring, so nobody reads the exclusion as an oversight.
\par
The cost is real and is accepted rather than argued away: every \texttt{dvc repro} of a \texttt{train} stage reruns \texttt{evaluate}, even when the model is identical. At the measured scale that is about three seconds for the evaluation, against the 43 seconds of the fit that preceded it, so it changes nothing about how the pipeline is used. If the evaluation ever becomes expensive -- the SC-06 masking sweep of \#39 is the candidate, at roughly 25 predict calls per variant -- option B becomes worth revisiting, and this entry is the record of what it would buy.},
  ai = {Alternative generation, Alternative assessment, Recommendation},
  response = {Accepted},
  assessment = {AI raised the trade as an explicit decision rather than letting a timestamp drift into an artefact unremarked, which is the failure mode this repository has already had once with a test run's validation summary. It named the alternative that preserves byte stability and the cost of that alternative, and it separated the reproducibility that is claimed from the reproducibility that is not, then wrote the test so that the exclusion of \texttt{model.\allowbreak{}json} is visible in the assertion instead of being silently absent.},
  aievidence = {Claude Code session on 2026-09-30 implementing issue \#37: two fits of each variant were compared byte for byte in a scratchpad before the test was written, and the payload files were confirmed identical while \texttt{model.\allowbreak{}json} differed.},
  evidence = {\href{https://github.com/mlops-2627q1-mds-upc/MLOPS_Recommenditos/blob/main/recommenditos/modeling/model.py}{\texttt{recommenditos/\allowbreak{}modeling/\allowbreak{}model.\allowbreak{}py}}, \texttt{fit\_\allowbreak{}variant} and \texttt{Model.\allowbreak{}save}; \texttt{tests/\allowbreak{}test\_\allowbreak{}model.\allowbreak{}py::test\_\allowbreak{}the\_\allowbreak{}payload\_\allowbreak{}of\_\allowbreak{}a\_\allowbreak{}refit\_\allowbreak{}is\_\allowbreak{}byte\_\allowbreak{}identical} and \texttt{::test\_\allowbreak{}training\_\allowbreak{}is\_\allowbreak{}reproducible}; \href{https://github.com/mlops-2627q1-mds-upc/MLOPS_Recommenditos/blob/main/docs/docs/requirements.md}{requirements} NFR-06; \hyperref[EDN-08]{EDN-08}; \hyperref[EDN-55]{EDN-55}; \href{https://github.com/mlops-2627q1-mds-upc/MLOPS_Recommenditos/issues/37}{issue \#37}.},
}
